\subsection{Classical logistic ODE}
            The classical baseline is the Verhulst logistic growth model
    \citep{verhulst1838notice}, fitted here as a shared-parameter nonlinear
    regression model. Unlike the neural models, it is fitted directly to the
    original response values $y_i^{[j]}$. The only normalization used by
    this baseline is the normalized time variable $\tau_i\in[0,1]$. Each
    trial has its own observed initial condition, fixed at the first
    measured value, but all trials share the same growth rate and carrying
    capacity:

    \begin{equation}
        \label{eqn:logisitc_ode}
        \frac{d y^{[j]}}{d\tau}
            =
            r y^{[j]}\left(1-\frac{y^{[j]}}{K}\right),
            \qquad
            y^{[j]}(0)=y_0^{[j]}=y^{[j]}(t_0).
    \end{equation}

    The exact solution used in the fit is

    \begin{equation}
        \label{eqn:sol_det_logistic}
        y^{[j]}(\tau)
            =
            \dfrac{K}{
                1 +
                \left(
                    \dfrac{K - y_0^{[j]}}{ y_0^{[j]} }
                \right)
                \exp(-r\tau)
            }.
    \end{equation}

        The unknown parameters of the classical fit are only $r$ and $K$; the
    initial values $y_0^{[j]}$ are taken from the data and are not
    optimized. The fit also requires positive observations, because the
    closed-form logistic curve and the ratio $(K-y_0^{[j]})/y_0^{[j]}$ are
    only meaningful for positive initial values.
    Missing observations, if present, are excluded through the observed
    index set $\Omega$.

        The fitted parameters must satisfy $r>0$ and $K>\max_j y_0^{[j]}$.
    Instead of optimizing constrained parameters directly, the
    implementation optimizes two unconstrained real variables
    $a, b \in\mathbb R$ and maps them to valid physical parameters with a
    $\operatorname{softplus}$ transformation:

    \begin{equation}
        r = \operatorname{softplus}(a)+10^{-6},
            \qquad
            K = K_{\min} +
            \operatorname{softplus}(b)+10^{-6},
    \end{equation}
    where
    $
        \displaystyle
        K_{\min}
     :=
        \max_{j \in J}
       \Big\{
                y_0^{[j]}
       \Big\}
            +10^{-3}
    $
    and
    $$
        \operatorname{softplus}(x)=\log(1+\exp(x)).
    $$

        The $\operatorname{softplus}$ function is a smooth differentiable
    approximation to the positive part function. It maps every real number
    to a positive value, behaves approximately like $\exp(x)$ for large
    negative $x$, and behaves approximately like $x$ for large positive
    $x$. This is useful for gradient-based fitting:
    Adam can update the unconstrained variables $a$ and $b$ freely, while
    the reported parameters $r$ and $K$ remain positive and the carrying
    capacity stays above every initial condition. The small constants
    $10^{-6}$ and $10^{-3}$ are numerical safeguards that avoid zero growth
    and avoid a carrying capacity equal to an initial value.

        The raw parameters are initialized by applying the inverse softplus
    map. In the
    implementation, this gives an initial growth rate near $3.0$ on the
    normalized time scale, and an initial carrying capacity slightly above the
    largest observed response. At each training epoch the code applies the
    softplus    transformation and evaluates the exact logistic solution for
    all trials and all observation times.

        The fitting criterion is the masked original-scale mean squared error
    \begin{equation}
        \label{eqn:likelihood_ode}
          \mathcal L_{\mathrm{logistic}}(r,K)
          =
            \frac{1}{n}
            \sum_{(i,j)\in\Omega}
            \left(y_i^{[j]}-\widehat y^{[j]}(\tau_i)\right)^2.
    \end{equation}

    This can also be interpreted as a maximum-likelihood fit after adding a
    Gaussian observation-error model on the original response scale:
    \begin{equation}
        y_i^{[j]}
        =
        \widehat y^{[j]}(\tau_i;r,K)+\varepsilon_i^{[j]},
        \qquad
        \varepsilon_i^{[j]}\stackrel{\mathrm{ind}}{\sim}
        \mathcal N(0,s^2).
    \end{equation}

        For fixed $s^2$, maximizing this Gaussian likelihood over $r$ and
    $K$ is equivalent to minimizing the sum of squared residuals.
    If $s^2$ is also unknown, profiling it out gives
    $$
        \widehat s^2(r,K)
        =
            \frac{1}{n}
            \sum_{(i,j)\in\Omega}
            \left(y_i^{[j]}-\widehat y^{[j]}(\tau_i;r,K)\right)^2,
    $$
and the maximizing values of $r$ and $K$ are again the nonlinear least-squares
solution. Therefore the implemented classical fit is a profiled Gaussian
maximum-likelihood estimate under an independent additive measurement-error
assumption. It is not a transition likelihood for a stochastic logistic process.

The code backpropagates through the closed-form expression above and applies an
Adam update \citep{kingma2015adam} to $a$ and $b$. No numerical ODE solver is
used for this classical baseline. This model was trained for 3000 Adam epochs
with learning rate $5\times10^{-2}$.

The inverse-problem procedure for this classical logistic model is therefore:

\begin{enumerate}
  \item Read the observed trial matrix and keep only entries in $\Omega$.
  \item Fix the initial condition of each trial at its observed first value
        $y_0^{[j]}$; these initial values are not optimized.
  \item Choose unconstrained optimization variables $a$ and $b$, and transform
        them into admissible physical parameters $r(a)>0$ and $K(b)>K_{\min}$
        with the softplus maps above.
  \item For the current values of $a$ and $b$, compute the exact logistic curve
        $\widehat y^{[j]}(\tau_i;r,K)$ for every trial and observation time.
  \item Evaluate the masked residual objective
        $\mathcal L_{\mathrm{logistic}}(r,K)$, using only observed entries.
  \item Use automatic differentiation through the closed-form solution to
        compute $\partial \mathcal L_{\mathrm{logistic}}/\partial a$ and
        $\partial \mathcal L_{\mathrm{logistic}}/\partial b$.
  \item Apply an Adam step to $a$ and $b$, then repeat until the fixed epoch
        budget is reached.
  \item Report $\widehat r=r(\widehat a)$ and
        $\widehat K=K(\widehat b)$, and evaluate fitted curves and residual
        metrics on the original data scale.
\end{enumerate}

In this baseline the inverse problem is low-dimensional: only two unknown
physical parameters are estimated. The use of gradient descent is an
implementation choice; mathematically, the estimator is the nonlinear
least-squares solution for the two-parameter logistic curve under fixed
trial-specific initial conditions.
The fitted values were
$$
  \widehat r = 7.705769,
  \qquad
  \widehat K = 47.860203.
$$
Since $\tau=t/168$ for these data, the growth rate in the original time scale
is approximately
$$
    \widehat r/168=0.045868
$$
per original time unit.